\section*{Chapitre \ref{ch:b1:lewis-resonance-vsepr} --- Schémas de Lewis, mésomérie et VSEPR}

\begin{solution}{exo:b1:lewis-resonance-vsepr:1}
\ce{H2O}~: deux liaisons \ce{O-H}, deux \omterm{def:b1:lewis-resonance-vsepr:lewis}{doublets non liants} sur O.
\ce{NH3}~: trois liaisons \ce{N-H}, un \omterm{def:b1:lewis-resonance-vsepr:lewis}{doublet non liant} sur N.
\ce{CO2}~: \ce{O=C=O}, deux \omterm{def:b1:lewis-resonance-vsepr:lewis}{doublets non liants} sur chaque O.
\ce{HCN}~: \ce{H-C#N}, un \omterm{def:b1:lewis-resonance-vsepr:lewis}{doublet non liant} sur N. \ce{CH2O}~: le carbone
est lié à deux H et doublement lié à O, deux \omterm{def:b1:lewis-resonance-vsepr:lewis}{doublets non liants} sur O.
Toutes les \omterm{def:b1:lewis-resonance-vsepr:formal-charge}{charges formelles} sont nulles dans ces cinq espèces.
\ce{NH4+}~: quatre liaisons \ce{N-H}, aucun \omterm{def:b1:lewis-resonance-vsepr:lewis}{doublet non liant}~; azote
$5 - 0 - 4 = +1$, la charge de l'ion.
\end{solution}

\begin{solution}{exo:b1:lewis-resonance-vsepr:2}
\ce{H2S}~: \ce{AX2E2}, coudée, angle inférieur à $109.5^\circ$ (mesuré~:
$92^\circ$). \ce{PCl3}~: \ce{AX3E}, pyramide trigonale, environ
$100^\circ$. \ce{BF3}~: \ce{AX3}, triangulaire plane, $120^\circ$.
\ce{SiH4}~: \ce{AX4}, tétraédrique, $109.5^\circ$. \ce{CO3^{2-}}~:
\ce{AX3} (la liaison double compte une fois), triangulaire plane,
$120^\circ$.
% ledger: geo:H2S.aHSH, geo:PCl3.aClPCl
\end{solution}

\begin{solution}{exo:b1:lewis-resonance-vsepr:3}
Apolaires~: \ce{CCl4} (\ce{AX4} tétraédrique), \ce{BF3} (\ce{AX3}),
\ce{CO2} (\ce{AX2} linéaire)~: les \omterm{def:b1:lewis-resonance-vsepr:dipole}{moments de liaison} se compensent.
Polaires~: \ce{CH2Cl2} (deux sortes de liaisons sur un tétraèdre),
\ce{NH3} (pyramidale), \ce{SO2} (coudée).
\end{solution}

\begin{solution}{exo:b1:lewis-resonance-vsepr:4}
\ce{H+} est l'\omterm{def:b1:lewis-resonance-vsepr:lewis-acid}{acide de Lewis} et \ce{H2O} la base~: la nouvelle liaison
\ce{O-H} de \ce{H3O+} est dative au moment de sa formation (ensuite les
trois liaisons \ce{O-H} sont identiques). \ce{Cu^{2+}} est l'acide et
chaque \ce{NH3} une base~: les quatre liaisons \ce{Cu-N} sont datives.
\end{solution}

\begin{solution}{exo:b1:lewis-resonance-vsepr:5}
\chemfig{CH_3-C(=[:60]O)-[:-60]\chemabove{O}{\scriptstyle\ominus}}
$\leftrightarrow$
\chemfig{CH_3-C(-[:60]\chemabove{O}{\scriptstyle\ominus})=[:-60]O}.
Les deux formules sont équivalentes~: les deux liaisons \ce{C-O} ont
l'indice 1.5 et la même longueur, intermédiaire entre la liaison double
et la liaison simple. Dans l'acide méthanoïque, les deux oxygènes
diffèrent (l'un porte le H), les formules ne sont pas équivalentes et
les liaisons gardent des longueurs distinctes, 120.2 (\ce{C=O}) et
\qty{134.3}{pm} (\ce{C-OH}).
\end{solution}

\begin{solution}{exo:b1:lewis-resonance-vsepr:6}
$N = 6 + 4 \times 6 + 2 = 32$ électrons. Avec quatre liaisons simples et
trois \omterm{def:b1:lewis-resonance-vsepr:lewis}{doublets non liants} sur chaque oxygène~: soufre $6 - 0 - 4 = +2$,
chaque oxygène $6 - 6 - 1 = -1$~; total $+2 - 4 = -2$. Géométrie
\ce{AX4}, tétraédrique. Avec deux liaisons \ce{S=O}, le choix des deux
oxygènes doublement liés parmi les quatre donne $\binom{4}{2} = 6$ formules.
\end{solution}

\begin{solution}{exo:b1:lewis-resonance-vsepr:7}
\ce{SF4}~: \ce{AX4E}, en bascule. \ce{ClF3}~: \ce{AX3E2}, en T.
\ce{XeF2}~: \ce{AX2E3}, linéaire. \ce{XeF4}~: \ce{AX4E2}, plan carré.
Dans une bipyramide trigonale, une position axiale a trois voisins à
$90^\circ$, une position équatoriale deux seulement~: les \omterm{def:b1:lewis-resonance-vsepr:lewis}{doublets non
liants}, encombrants, se placent là où ils subissent le moins de
répulsions à $90^\circ$.
\end{solution}

\begin{solution}{exo:b1:lewis-resonance-vsepr:8}
Le soufre et le phosphore sont plus gros et moins électronégatifs que
l'oxygène et l'azote~: les \omterm{def:b1:lewis-resonance-vsepr:lewis}{doublets liants} sont plus loin de l'atome
central et attirés vers l'hydrogène, ils se repoussent donc moins, et
les \omterm{def:b1:lewis-resonance-vsepr:lewis}{doublets non liants} resserrent les liaisons vers $90^\circ$.
\end{solution}

\begin{solution}{exo:b1:lewis-resonance-vsepr:9}
\ce{CH3-C#N}~: le carbone du \ce{CH3} est sp$^3$, le carbone du nitrile
sp, l'azote sp~; 5 liaisons $\sigma$ (trois \ce{C-H}, \ce{C-C}, une dans
\ce{C#N}) et 2 liaisons $\pi$. \ce{CH2=CH-CH3}~: les deux carbones de
l'alcène sont sp$^2$, le carbone du méthyle sp$^3$~; 8 liaisons $\sigma$
(six \ce{C-H}, deux \ce{C-C}) et 1 liaison $\pi$.
\end{solution}

\begin{solution}{exo:b1:lewis-resonance-vsepr:10}
Les deux molécules sont pyramidales, avec un \omterm{def:b1:lewis-resonance-vsepr:lewis}{doublet non liant} sur
l'azote. Le \omterm{def:b1:lewis-resonance-vsepr:lewis}{doublet non liant} apporte une contribution orientée (du
négatif vers le positif) du doublet vers le noyau. Dans \ce{NH3}, l'azote
est l'extrémité négative de chaque liaison~: les \omterm{def:b1:lewis-resonance-vsepr:dipole}{moments de liaison} sont
orientés de N vers les hydrogènes, dans le même sens que la contribution
du doublet, et ils s'ajoutent. Dans \ce{NF3}, c'est le fluor qui est
l'extrémité négative~: les \omterm{def:b1:lewis-resonance-vsepr:dipole}{moments de liaison} sont orientés des fluors
vers N, à l'opposé de la contribution du doublet, et les deux se
compensent presque.
\end{solution}

\begin{solution}{exo:b1:lewis-resonance-vsepr:11}
$\mu_b = \mu/(2\cos(\theta/2)) = 1.63/(2\cos 59.75^\circ) =
\qty{1.62}{\debye} = \qty{5.40e-30}{C.m}$. D'où $\delta = \mu_b/(e\,d) =
5.40 \times 10^{-30}/(1.602 \times 10^{-19} \times 143.2 \times 10^{-12})
= 0.24$.
\end{solution}

\begin{solution}{exo:b1:lewis-resonance-vsepr:12}
\ce{N=N=O} avec les \omterm{def:b1:lewis-resonance-vsepr:formal-charge}{charges formelles} $-1$, $+1$, $0$~; \ce{N#N-O} avec
$0$, $+1$, $-1$. La liaison \ce{N-N} (\qty{112.8}{pm}) est proche de la
liaison triple de \ce{N2} (\qty{109.8}{pm})~: la formule \ce{N#N-O} a le
plus de poids, et elle place la charge négative sur l'oxygène, l'atome le
plus électronégatif. La liaison \ce{N-O} est cependant plus courte
qu'une liaison simple~: l'autre formule contribue aussi.
\end{solution}

\begin{solution}{pb:b1:lewis-resonance-vsepr:1}
\textbf{1.} \ce{N2} 10, \ce{NO} 11, \ce{NO2} 17, \ce{N2O} 16, \ce{O3} 18,
\ce{HNO3} 24.
\textbf{2.} \ce{N#N}, avec un \omterm{def:b1:lewis-resonance-vsepr:lewis}{doublet non liant} sur chaque azote.
\textbf{3.} \ce{N=O}, avec deux \omterm{def:b1:lewis-resonance-vsepr:lewis}{doublets non liants} sur l'oxygène, un
\omterm{def:b1:lewis-resonance-vsepr:lewis}{doublet non liant} et un \omterm{def:b1:quantum-numbers:unpaired}{électron célibataire} sur l'azote. Avec 11
électrons, nombre impair, un atome a toujours un compte impair~: l'azote
est entouré de 7 électrons.
\textbf{4.} \ce{N=N=O}~: $-1$, $+1$, $0$~; \ce{N#N-O}~: $0$, $+1$, $-1$.
\textbf{5.} \ce{O=O-O}~: oxygène central $+1$ (un \omterm{def:b1:lewis-resonance-vsepr:lewis}{doublet non liant}, trois
liaisons), oxygène terminal lié par une liaison simple $-1$, l'autre 0.
\textbf{6.} Comme dans l'\cref{ex:b1:lewis-resonance-vsepr:hno3}~: azote
$+1$, l'oxygène portant une liaison simple et pas d'hydrogène $-1$, les
autres 0.
\textbf{7.} \ce{NO} et \ce{NO2} (11 et 17 électrons).
\textbf{8.} Deux formules équivalentes, la liaison double d'un côté ou de
l'autre~: indice 1.5 pour chaque liaison.
\textbf{9.} Oui~: \qty{127.8}{pm} est compris entre la liaison double
(\qty{120.8}{pm}) et la liaison simple (\qty{147.5}{pm}).
\textbf{10.} La liaison double tour à tour sur chacun des trois oxygènes~;
chaque liaison \ce{N-O} est double dans une formule sur trois~: indice
$4/3$.
\textbf{11.} La liaison longue, \qty{140.6}{pm}, est \ce{N-OH}, une
liaison simple. Les deux oxygènes sans hydrogène se partagent une liaison
$\pi$ à travers deux \omterm{def:b1:lewis-resonance-vsepr:resonance}{formules mésomères} équivalentes~: deux liaisons
d'indice 1.5, 121.1 et \qty{119.9}{pm}.
\textbf{12.} L'azote porte un \omterm{def:b1:quantum-numbers:unpaired}{électron célibataire} au lieu d'un \omterm{def:b1:lewis-resonance-vsepr:lewis}{doublet
non liant}~: il repousse les \omterm{def:b1:lewis-resonance-vsepr:lewis}{doublets liants} moins qu'un doublet ne le
ferait, et les liaisons s'ouvrent au-delà de $120^\circ$.
\textbf{13.} \ce{N#N-O}, avec la charge négative sur l'oxygène.
\textbf{14.} \ce{O3}~: \ce{AX2E}, coudée. \ce{N2O}~: \ce{AX2} (N central),
linéaire. \ce{NO3-}~: \ce{AX3}, triangulaire plane. Azote de \ce{HNO3}~:
\ce{AX3}, triangulaire plane.
\textbf{15.} Le \omterm{def:b1:lewis-resonance-vsepr:lewis}{doublet non liant} de l'oxygène central repousse les
\omterm{def:b1:lewis-resonance-vsepr:lewis}{doublets liants} plus qu'ils ne se repoussent entre eux.
\textbf{16.} L'ozone est coudé et son atome central diffère des atomes
terminaux (\omterm{def:b1:lewis-resonance-vsepr:formal-charge}{charge formelle} positive au centre, charge négative partagée
par les extrémités)~: un faible \omterm{def:b1:lewis-resonance-vsepr:dipole}{moment dipolaire} selon la bissectrice.
\ce{N2O} est linéaire mais ses deux extrémités sont des atomes
différents~: un faible \omterm{def:b1:lewis-resonance-vsepr:dipole}{moment dipolaire}. \ce{N2} est symétrique~: aucun.
\textbf{17.} \ce{CO2} est linéaire (\ce{AX2})~: ses \omterm{def:b1:lewis-resonance-vsepr:dipole}{moments de liaison} se
compensent. \ce{SO2} est coudé (\ce{AX2E})~: ils s'ajoutent.
\textbf{18.} \ce{AX3E}~: le \omterm{def:b1:lewis-resonance-vsepr:lewis}{doublet non liant} referme les angles
en dessous de $109.5^\circ$~; mesuré~: $106.7^\circ$.
\textbf{19.} \ce{AX2E2}~: deux \omterm{def:b1:lewis-resonance-vsepr:lewis}{doublets non liants} le referment davantage~;
$104.5^\circ$.
\textbf{20.} $\mu = 2\mu_{OH}\cos(\theta/2)$
(\cref{prop:b1:lewis-resonance-vsepr:dipole-sum}).
\textbf{21.} $\mu_{OH} = 1.857/(2\cos 52.24^\circ) = 1.857/1.2246 =
\qty{1.52}{\debye}$.
\textbf{22.} $1.516 \times 3.336 \times 10^{-30} = \qty{5.06e-30}{C.m}$.
\textbf{23.} $e\,d = 1.602 \times 10^{-19} \times 95.8 \times 10^{-12} =
\qty{1.535e-29}{C.m}$, soit \qty{4.60}{\debye}.
\textbf{24.} $\delta = 5.06 \times 10^{-30}/1.535 \times 10^{-29} =
\textbf{0.33}$~: la liaison \ce{O-H} de l'eau porte le tiers de la
charge que porterait une liaison entièrement ionique.
\end{solution}
